\section{Discussion}
\label{sec:discussion}
Spatial CV is now standard for validating ecological and environmental maps, because random CV can severely understate prediction error when the target lies far from the data \citep{roberts2017cross,ploton2020spatial,meyer2022global}. Nearest-neighbour distance matching addresses this by making validation resemble prediction geographically \citep{mila2022nndm,linnenbrink2024knndm}. Our results clarify what such matching guarantees about prediction error and what it leaves undetermined. Three findings stand out. First, for Gaussian prediction from a single neighbour with known covariance, correlation-weighted matching is exactly the right criterion: $VJ_2$ equals the expected CV bias. Second, once prediction uses several observations, the information they supply jointly can bias CV even when nearest distances match perfectly. Third, correlation weighting improved scheme selection substantially when covariance was known, but by only \NestedReduction\% when covariance was estimated, and almost entirely in one setting.

\subsection{Guarantees and limits of nearest-neighbour matching}
The one-neighbour identity, $O_1=VJ_2$, explains why distances matter only through the correlation they imply. A shift in nearest distance changes prediction error when it changes the information a neighbour supplies. Where correlation decays rapidly, a small shift can remove most of that information. Where responses are weakly correlated at both validation and target distances, a large shift may have little effect. The Wasserstein distance $W$ weights every unit of distance equally and therefore cannot distinguish these situations, whereas $J_2$ weights CDF differences by the rate at which squared correlation declines (equation~\eqref{eq:weighting}). NNDM already incorporates dependence through an autocorrelation range \citep{mila2022nndm}; $J_2$ makes that dependence explicit and connects it directly to MSE.

The decomposition in equation~\eqref{eq:remainder} identifies what any nearest-neighbour summary leaves out. The remainder compares how much the other training observations reduce error during validation and at the target. It is small when the nearest observation largely screens the others from the prediction location, a property of kriging studied by \citet{stein2002screening}. Screening depends on the covariance model, the arrangement of observations and measurement error \citep{stein1999interpolation,stein2002screening}, so it cannot be assumed in general. The leave-one-out counterexample shows that the remainder alone can produce nonzero expected bias under a correctly specified model with perfectly matched nearest distances. Perfect distance matching is therefore not sufficient for unbiased validation. Nor can weights based only on nearest distances repair this: in the counterexample, such weights are all equal and leave the CV average unchanged.

The bound in equation~\eqref{eq:certificate} quantifies the gap. When its interval excludes zero, it determines the direction of expected CV bias under the stated covariance model. It contained every exact MSE difference in our simulations, but often included zero. This is the expected price of a guarantee that holds without further assumptions about how training observations are arranged.

\subsection{Evidence for correlation-weighted scheme selection}
The simulation evidence separates two targets of inference. With known covariance, selecting by $|J_2|$ rather than $W$ reduced mean absolute expected bias by \BiasReduction\%, and $J_2$ identified the direction of bias in \SignJTwo\% of evaluations with appreciable bias. This is the quantity that the theory addresses. With covariance estimated within each fold, the relevant quantity is the error of the observed CV score for the model fitted to that dataset, which also varies with the particular responses observed \citep{bates2024crossvalidation}. Here the reduction was \NestedReduction\%, concentrated in a setting with uniform sampling, a target extending beyond the sampled region, rapidly decaying correlation and low noise. Across the other 23 settings the average gain was negligible, and the ten-fold check gave an interval that included zero (supplementary material, Section~\ref*{supp-sec:tenfold}).

Two mechanisms explain this attenuation. First, estimated covariance degrades the contrast itself: sign agreement fell from \SensIsoTrue\% to \SensIsoEst\% under isotropic dependence, and fitted variograms sometimes placed the signal fraction or correlation scale at extreme values. Estimating covariance parameters also adds prediction error that plug-in calculations do not capture \citep{zimmerman1992mspe}. Second, observed CV scores vary between response datasets. The MAE of the scores selected by $W$ was \NestedMAEW, several times the mean absolute expected bias of \BiasW\ in the known-covariance study, although the two designs differ. When this variation is large relative to the differences in expected bias between candidate schemes, choosing the scheme with the best expected properties yields little gain for a single dataset. Comparisons between CV procedures are subject to substantial sampling variability more generally \citep{lei2020cvc,bates2024crossvalidation}.

The Jura application illustrates the practical consequence. Both diagnostics selected kNNDM for the original 100 validation sites, yet random CV was closer to the observed target MSE. With the northernmost sites as target, the selected block schemes were closer than random CV. The scheme that best matches nearest-neighbour summaries need not give the most accurate estimate for a particular survey, and the empirical comparison depends on the target. Because the application uses one survey, it cannot separate the additional-neighbour term from variation between response datasets.

\subsection{Spatial proximity is not data leakage}
The prediction model determines which spatial information matters. In Section~\ref{sec:sim-learners}, the bias of random CV for a random forest almost vanished as the predictor became informative, although response correlation and $J_2$ remained fixed. Response correlation describes dependence before accounting for predictors. Residual dependence, the model's use of predictors and whether target predictor values are represented in training all influence CV bias. Distance-based error profiles and area-of-applicability methods address the last two directly \citep{brenning2023profiles,meyer2021aoa}.

These diagnostics should therefore not be read as measures of data leakage. Leakage means using information that will not be available when the model is deployed \citep{kaufman2012leakage}. Nearby observations may be entirely legitimate at deployment; whether they are depends on the prediction task. The contrasts compare the spatial conditions of validation and prediction for a declared target. A contrast near zero does not certify that a CV estimate is accurate. It can arise from similar distance distributions, cancellation between positive and negative CDF differences, or weak correlation at all relevant distances. Like any single index, it could also become a target to optimise rather than evidence to interpret.

\subsection{Recommendations for practice}
The results support four recommendations for spatial validation.
\begin{enumerate}
\item \emph{Declare the prediction target.} CV bias is defined relative to target locations and their weights. Different targets in the same survey favoured different schemes in the Jura application.
\item \emph{Report the distance distributions together with correlation-weighted summaries.} The CDFs and $W$ describe the geographic mismatch. Comparing $|J_2|$ with $U_2$ reveals whether a small contrast arises from cancellation.
\item \emph{Treat diagnostics as conditional on a model.} $J_2$ and the bound rely on a Gaussian covariance model. Prediction error profiles, applicability assessments and residual analyses address what response correlation cannot \citep{brenning2023profiles,meyer2021aoa}.
\item \emph{Use probability sampling to estimate map accuracy.} For the error of a fixed map over a region, an independent probability sample supports design-based estimation without model assumptions \citep{wadoux2021spatial}.
\end{enumerate}

\subsection{Limitations and future directions}
The exact identities and bounds assume Gaussian responses with known mean and covariance. For non-Gaussian responses, such as presence--absence data, the relationship between correlation and MSE changes, and comparable results would require model-based extensions \citep{diggle1998model}. The simulations used exponential and Mat\'ern-$3/2$ covariance, a known mean and a fixed grid of designs, and overall MAE weights these settings equally. Inserting parameter estimates into the identities does not account for estimation uncertainty \citep{zimmerman1992mspe}, and inference after data-dependent scheme selection remains open.
The Jura analysis checked the covariance-fitting objective with a profile search over correlation scales, which left its findings unchanged (supplementary material, Section~\ref*{supp-sec:jura-fitting}). Rerunning all three estimated-covariance studies with the same search changed 5.5--9.4\% of variogram fits by more than 1\% in correlation scale, yet the selection results were essentially unaltered: the five-fold MAE reduction was 2.4\% (interval excluding zero), the ten-fold reduction 1.5\% (interval including zero), and the benefit remained concentrated in the same setting. Sign agreement with estimated covariance changed by 0.3 percentage points (supplementary material, Section~\ref*{supp-sec:profile-refits}). The reported gains are therefore not an artefact of the fitting procedure, although they remain specific to the settings studied.

Several extensions follow directly from the decomposition. Covariance-based corrections modify the error estimator rather than choosing among schemes \citep{rabinowicz2022correlated}, and importance weighting adjusts for differences between sampled and target distributions \citep{sugiyama2007covariate,debruin2022clustered,brenning2026targetweighted}. Comparing these approaches with diagnostic scheme selection is an open problem. Any correction that improves on nearest-neighbour matching must represent the joint information supplied by several neighbours and the uncertainty in estimated dependence. Equation~\eqref{eq:remainder} specifies that term exactly.

Distance matching makes spatial validation resemble the intended prediction geographically. Our results show that this resemblance determines CV bias exactly only when prediction depends on a single neighbour and the covariance is known. Beyond that setting, credible accuracy assessment requires the prediction target, the prediction model and independent validation data, as well as the geometry of the sample.
